% !TEX program = lualatex
% !TeX encoding = UTF-8
\PassOptionsToPackage{unicode}{hyperref}
\PassOptionsToPackage{bookmarksnumbered,bookmarksopen}{hyperref}
\documentclass[12pt,a4paper]{article}

% ============================================================
% 한국어 설정 (LuaLaTeX + luatexja)
% ============================================================
\usepackage{luatexja}
\usepackage{luatexja-fontspec}
\usepackage[english]{babel}   % 영어용 babel (필요시)

% 한국어 고딕체 (sans) – Noto Sans KR
\setsansjfont{Noto Sans KR}[
UprightFont = * Regular,
BoldFont = * Bold,
Script = Hangul,
Language = Korean
]

% 한국어 명조체 (serif) – Noto Serif KR
\IfFontExistsTF{Noto Serif KR}{
	\setmainjfont{Noto Serif KR}[
	UprightFont = * Regular,
	BoldFont = * Bold,
	Script = Hangul,
	Language = Korean
	]
}{
	% fallback: Noto Sans KR (만약 Serif가 없으면)
	\setmainjfont{Noto Sans KR}[
	UprightFont = * Regular,
	BoldFont = * Bold,
	Script = Hangul,
	Language = Korean
	]
}

% 고정폭 폰트 – 라틴 문자는 Latin Modern Mono, 한글은 Noto Sans KR
\setmonojfont{Noto Sans KR}[
UprightFont = * Regular,
BoldFont = * Bold,
Script = Hangul,
Language = Korean
]

% 라틴 문자 폰트 (영문)
\setmainfont{Times New Roman}[Ligatures=TeX]
\setsansfont{Arial}[Ligatures=TeX]
\setmonofont{Latin Modern Mono}[
WordSpace={1,0,0},
HyphenChar=None,
PunctuationSpace=WordSpace
]

% ============================================================
% 패키지 (원본과 동일)
% ============================================================
\usepackage{amsmath,amssymb,amsthm,mathtools}
\usepackage{bm}
\usepackage{braket}
\usepackage{siunitx}
\usepackage{graphicx}
\usepackage{tikz}
\usepackage{tikz-3dplot}
\usepackage{pgfplots}
\pgfplotsset{compat=1.18}
\usetikzlibrary{arrows.meta,positioning,shapes,calc,angles,quotes,patterns,3d}

\usepackage{booktabs}
\usepackage{listings}
\usepackage{xcolor}
\usepackage{algorithm}
\usepackage{algorithmic}
\usepackage{multicol}
\usepackage{enumitem}
\usepackage{tcolorbox}
\usepackage{hyperref}
\usepackage{cleveref}

% 기하학 설정 (margin=2.5cm 제거, 기본값 사용)
\usepackage{geometry}
% \geometry{margin=2.5cm}  % 주석 처리 또는 삭제

% 정리 환경 (한국어)
\newtheorem{theorem}{정리}
\newtheorem{lemma}{보조정리}
\newtheorem{corollary}{계}
\newtheorem{definition}{정의}
\newtheorem{axiom}{공리}
\newtheorem{proposition}{명제}
\newtheorem{remark}{주석}
\newtheorem{assumption}{가정}

% 연산자
\DeclareMathOperator{\Tr}{Tr}
\DeclareMathOperator{\Diff}{Diff}
\DeclareMathOperator{\Aut}{Aut}
\DeclareMathOperator{\Vol}{Vol}
\DeclareMathOperator{\Ric}{Ric}
\DeclareMathOperator{\Riem}{Riem}
\DeclareMathOperator{\Cov}{Cov}
\DeclareMathOperator{\supp}{supp}
\DeclareMathOperator{\diag}{diag}
\DeclareMathOperator{\sign}{sign}
\DeclareMathOperator{\dimension}{dim}
\DeclareMathOperator{\Div}{Div}
\DeclareMathOperator{\Grad}{Grad}
\DeclareMathOperator{\Hess}{Hess}
\DeclareMathOperator{\Ricci}{Ricci}
\DeclareMathOperator{\Scalar}{Scalar}

\DeclareRobustCommand{\alD}{\texorpdfstring{\ensuremath{\alpha_{\mathrm{D}}}}{alpha_D}}
\DeclareRobustCommand{\beI}{\texorpdfstring{\ensuremath{\beta_{\mathrm{I}}}}{beta_I}}
\DeclareRobustCommand{\gaR}{\texorpdfstring{\ensuremath{\gamma_{\mathrm{R}}}}{gamma_R}}
\DeclareRobustCommand{\UCS}{\texorpdfstring{\ensuremath{\mathcal{M}_{\mathrm{UCS}}}}{M_UCS}}

% YUCT 고유 명령어
\newcommand{\eff}{\mathrm{eff}}
\newcommand{\coord}{\mathrm{coord}}
\newcommand{\Yak}{\mathrm{Yak}}
\newcommand{\Dict}{\mathcal{D}}
\newcommand{\Mdict}{\mathcal{M}_{\Dict}}
\newcommand{\Ke}{K_{\eff}}
\newcommand{\Kemax}{K_{\eff,\max}}
\newcommand{\Keobs}{K_{\eff,\mathrm{obs}}}
\newcommand{\Keexp}{K_{\eff,\mathrm{exp}}}
\newcommand{\Kec}{K_{\eff,c}}
\newcommand{\Kcrit}{K_{\mathrm{crit}}}
\newcommand{\Keff}{K_{\mathrm{eff}}}

% PDF 메타데이터
\hypersetup{
	pdftitle={부록 K: 종교적 실천을 YPSDC 프로토콜로 보기 – 야쿠셰프 이론에 의한 형식화},
	pdfauthor={Alexey V. Yakushev},
	pdfsubject={YUCT, YPSDC, 종교적 실천, 조정 이론, 비잔틴 음향, 검증 프로그램},
	pdfkeywords={YUCT, YPSDC, 종교적 실천, 조정 이론, 비잔틴 음향, 검증, 학생 연구, 학제간}
}

\begin{document}
	
	\begin{titlepage}
		\begin{center}
			\vspace*{0cm}
			
			\Huge\textbf{부록 K: 종교적 실천을 YPSDC 프로토콜로 보기: 야쿠셰프 이론에 의한 형식화}
			
			\vspace{1cm}
			
			\small\textit{비잔틴 음향에서 실험적 검증 프로그램까지의 통일적 프레임워크}
			
			\vspace{0.5cm}
			
			\small\textbf{Alexey V. Yakushev}
			
			\vspace{0cm}
			
			YUCT 이론: \url{https://yuct.org/}\\
			YPSDC 프로토콜: \url{https://ypsdc.com/}\\
			
			\vspace{2cm}
			\textbf{DOI: \url{https://doi.org/10.5281/zenodo.18444598}}\\
			
			\vspace{1cm}
			
			\vspace{0cm}
			\large 
			본 연구의 축약판은 이전에 출판되었으며 다음에서 확인할 수 있습니다: \\ \url{https://doi.org/10.5281/zenodo.18362308}\\
			\vspace{1cm}
			2026년 3월 \\ \large V.2.1.
			\newpage
			\begin{abstract}
				\noindent
				본 문서는 야쿠셰프 통합 조정 이론(YUCT)과 YPSDC(야쿠셰프 동기 분산 조정 프로토콜)의 관점에서 종교적 실천을 포괄적으로 형식화합니다. 제1부에서는 종교 의식이 사전 분배된 사전을 사용하는 고도로 최적화된 조정 프로토콜로서 기능함을 보여주고, 비잔틴 교회의 음향이 조정 효율을 위한 건축적 최적화의 전형적인 예임을 입증합니다. 제2부에서는 YUCT 프레임워크의 극단적인 통합력을 개괄하고, 실험적 검증을 위한 상세한 연구 프로그램을 제시하며, 학생 연구를 과학적 검증의 초석으로 강조합니다. 이 프레임워크는 측정 가능한 지표($K_{\mathrm{eff}}$), 검증 가능한 예측, 물리학·생물학에서 사회학·종교학에 이르는 학제간 응용을 제공합니다.
			\end{abstract}
			
			\vspace{0cm}
			
			\noindent
			\small\textbf{키워드:} YUCT, YPSDC, 종교적 실천, 조정 이론, 비잔틴 음향, 검증 프로그램, 학생 연구, 학제간, $K_{\mathrm{eff}}$ 측정
			
			\vspace{0.5cm}
			
			\noindent
			\textcopyright 2026 Yakushev Research. All rights reserved.
		\end{center}
	\end{titlepage}
	
	\tableofcontents
	
	\newpage
	
	\part{종교적 실천을 YPSDC 프로토콜로 보기: 야쿠셰프 이론에 의한 형식화}
	
	\section{서론: 종교적 조정의 모델로서의 YPSDC}
	
	야쿠셰프 조정 이론과 그 YPSDC(야쿠셰프 동기 분산 조정 프로토콜)는 종교적 실천을 사전 분배된 사전의 원리를 이용한 고도로 최적화된 조정 시스템으로 보는 새로운 시각을 제공한다.
	
	YPSDC는 두 단계에 기반한다:
	\begin{enumerate}
		\item \textbf{오프라인 단계}: 참가자들 사이에 사전(정전 텍스트, 기도, 성가, 의식)을 분배.
		\item \textbf{온라인 단계}: 복잡한 행동의 동기 활성화를 위한 짧은 코드(예: 기도의 시작)를 전송.
	\end{enumerate}
	
	종교 의식은 이 체계에 완벽하게 부합한다: 신자들은 미리 기도와 성가를 배우고(사전), 예배 중에 짧은 신호(기도의 시작, 사제의 외침)를 받아 동기적으로 행동한다.
	
	\section{종교적 조정의 YPSDC에 의한 수학적 형식화}
	
	\subsection{종교 시스템을 YPSDC 프로토콜로 정의}
	
	종교 조정 시스템을 튜플로 정의한다:
	\[
	R_{\mathrm{YPSDC}} = (A, P, D_{\mathrm{relig}}, K, C, T, R)
	\]
	여기서,
	\begin{itemize}
		\item $A = \{a_1, a_2, \ldots, a_N\}$ — 종교적 행동의 집합(기도, 성가, 의식적 동작)
		\item $P = \{P_1, P_2, \ldots, P_M\}$ — 참가자(신자, 성직자)
		\item $D_{\mathrm{relig}} = \{(\kappa_i, a_i)\}$ — 종교 사전. $\kappa_i$는 짧은 코드, $a_i$는 해당 행동
		\item $K$ — 전송되는 코드의 집합
		\item $C$ — 전달 채널(음향, 시각)
		\item $T$ — 시간적 제약(전례 주기)
		\item $R$ — 공명 인자(음향적, 사회적)
	\end{itemize}
	
	\subsection{종교적 조정의 효율}
	
	종교 실천의 조정 효율:
	\[
	K_{\mathrm{eff}}^{\mathrm{relig}} = \frac{H(A_{\mathrm{full}})}{H(\kappa_{\mathrm{code}})} \cdot R_{\mathrm{arch}} \cdot R_{\mathrm{soc}}
	\]
	여기서,
	\begin{itemize}
		\item $H(A_{\mathrm{full}})$ — 종교적 행동의 완전한 기술의 엔트로피(텍스트+선율+의식)
		\item $H(\kappa_{\mathrm{code}})$ — 활성화 코드의 엔트로피
		\item $R_{\mathrm{arch}}$ — 음향-건축적 공명 인자
		\item $R_{\mathrm{soc}}$ — 사회적 공명
	\end{itemize}
	
	\textbf{예: 주기도문}
	\begin{itemize}
		\item 완전한 기술: 약 1000비트(텍스트+선율+의식적 동작)
		\item 활성화 코드: 10~20비트("우리 아버지" 또는 첫 몇 마디)
		\item $K_{\mathrm{eff}}^{\mathrm{relig}} \approx 50$–$100$
	\end{itemize}
	
	\section{비잔틴 음향을 YPSDC의 최적화로}
	
	\subsection{사전 분배의 최적화}
	
	비잔틴 교회는 종교 사전의 분배와 활성화를 최적화했다:
	
	\begin{enumerate}
		\item \textbf{음향 증폭}: 돔 구조는 주요 성가의 음정(110 Hz — 남성 성가, 220 Hz — 여성 성가)에 맞는 공명 주파수를 창출했다.
		\item \textbf{시간적 동기화}: 2~4초의 잔향은 구절의 원활한 중첩을 보장하여 "지속적인 기도"의 효과를 창출했다.
		\item \textbf{공간적 배치}: 성가대를 특정 위치(성가대석, 소레아)에 배치하여 음향적 연결성을 최대화했다.
	\end{enumerate}
	
	\subsection{아야 소피아의 모드 분석}
	
	대성당의 공명 주파수:
	\[
	f_{n,m,l} = \frac{c}{2} \sqrt{\left(\frac{n}{31}\right)^2 + \left(\frac{m}{56}\right)^2 + \left(\frac{l}{55}\right)^2} \, \mathrm{m}
	\]
	여기서 주목할 만한 모드:
	\begin{itemize}
		\item 110 Hz (남성 성가의 기본음)
		\item 8–12 Hz (뇌의 알파 리듬 대역, 바이노럴 비트로 달성)
	\end{itemize}
	
	\subsection{음향 증폭 계수}
	\[
	R_{\mathrm{acoust}} = \frac{\int_V |p_{\mathrm{res}}(\mathbf{r})|^2 dV}{\int_V |p_{\mathrm{free}}(\mathbf{r})|^2 dV} \approx 200\text{--}500
	\]
	
	\section{실험적 측정 프로토콜}
	
	\subsection{프로토콜 A: YPSDC 동기화}
	
	\begin{enumerate}
		\item \textbf{준비 단계}: 참가자는 기도 사전(텍스트, 선율)을 학습한다.
		\item \textbf{활성화 단계}: 지도자가 활성화 코드(기도의 시작)를 발성한다.
		\item \textbf{측정}:
		\begin{itemize}
			\item 동기화 시간: $\tau_{\mathrm{sync}}$
			\item 재현 정확도: $F_{\mathrm{acc}}$
			\item 동기성: $\Phi_{ij} = \langle e^{i[\phi_i(t) - \phi_j(t)]} \rangle$
		\end{itemize}
		\item \textbf{메트릭}:
		\[
		K_{\mathrm{eff}}^{\mathrm{YPSDC}} = \frac{\tau_{\mathrm{without\,dict}}}{\tau_{\mathrm{with\,dict}}} \cdot \frac{1}{N} \sum_{i<j} \Phi_{ij}
		\]
	\end{enumerate}
	
	\subsection{프로토콜 B: 건축적 최적화}
	
	\begin{enumerate}
		\item \textbf{환경 비교}:
		\begin{itemize}
			\item 최적 음향의 사원
			\item 음향적으로 "데드"한 방
			\item 개방 공간
		\end{itemize}
		\item \textbf{매개변수}:
		\begin{itemize}
			\item 잔향 시간 $\tau_{60}$
			\item 음성 전달 지수 STI
			\item 명료도 계수 $C_{80}$
		\end{itemize}
		\item \textbf{의존 관계}:
		\[
		K_{\mathrm{eff}}^{\mathrm{arch}} = K_0 \cdot \left(1 + \alpha \frac{V}{V_0}\right) \cdot e^{-\tau/\tau_0} \cdot \mathrm{STI}
		\]
	\end{enumerate}
	
	\subsection{사례 연구: 이슬람 기도를 YPSDC 프로토콜로 – 동기화의 실증적 측정}
	\label{subsec:islamic_prayer}
	
	이슬람 기도(살라트)는 YPSDC 프로토콜의 매우 명확하고 전 세계적으로 접근 가능한 예를 제공한다. 하루 다섯 번의 기도는 태양의 위치에 따라 시간이 정해져 자연적인 행성 규모의 일정을 만든다. 기도 의식 자체는 어린 시절부터 모든 독실한 무슬림이 알고 있는 낭송과 신체 동작의 고도로 구조화된 연속이며, 이는 오프라인에서 분배된 사전의 완벽한 예이다.
	
	\subsubsection{시간과 기술에 의한 전 세계적 동기화}
	정확한 기도 시간은 지리적 위치에 따라 다르며 천문학적 공식을 사용하여 계산된다. 오늘날 이러한 시간은 모바일 앱, 웹사이트, 지역 모스크의 발표를 통해 전달된다. 아잔(예배의 부름)은 온라인 인덱스로 기능한다: 짧은 청각 신호(약 2~5분)가 수백만 명의 사람들에게 동시에 기도의 전체 시퀀스를 활성화한다.
	
	우리는 다음과 같은 전 세계적 규모의 측정을 제안한다:
	\begin{itemize}
		\item \textbf{데이터 소스}: 인기 기도 앱(예: Muslim Pro, Athan)에서 기도 알림 후 앱을 열거나 기도를 완료로 표시한 사용자에 대한 익명화된 데이터. 기도 관련 해시태그(\#Fajr, \#Dhuhr 등)를 포함한 소셜 미디어 활동(트윗, 게시물)의 타임스탬프 및 위치 정보.
		\item \textbf{분석}: 각 기도 시간에 대해 활동 강도의 시계열을 구성한다. 예상 기도 시작 시간(천문학적 계산에서)과 관측된 활동 피크 사이의 상호 상관을 계산한다. 피크의 날카로움(반치폭)은 동기화 정밀도 $\Delta t$를 반영한다.
		\item \textbf{조정 효율}: 기도 의식의 정보량은 낭송되는 텍스트의 길이와 복잡성 및 다양한 동작의 수(예: 한 번의 기도당 약 $10^4$비트)에서 추정할 수 있다. 인덱스(아잔)는 수백 비트만을 가진다. 따라서 단순한 $\Keff = \frac{\text{의식 정보}}{\text{인덱스 정보}} \approx 10$–$10^2$이다. 그러나 전 세계적 규모로 인해 동일한 인덱스가 수십억 명의 사람들을 동시에 조정하므로 실효적인 $\Keff$는 훨씬 더 클 수 있다.
	\end{itemize}
	
	\subsubsection{공동 기도에서의 국소적 생리적 동기화}
	많은 연구가 집단 의식이 참가자들 사이에 생리적 동기화를 유도함을 보여준다. 이슬람 기도에 대해 왕립학회의 연구 \cite{RoyalSociety2024}는 동기화된 동작 중에 유의미한 심박 변이 동기화와 호흡 동조를 발견했다. 우리는 이러한 측정을 확장하여 YPSDC 예측을 명시적으로 검증할 것을 제안한다:
	\begin{itemize}
		\item \textbf{참가자}: 모스크에서 공동 기도를 수행하는 10~50명의 자원자 그룹으로, 휴대용 ECG/HRV 센서와 가속도계를 착용한다.
		\item \textbf{조건}: 음향 환경(잔향 시간, 스피커 유무), 이맘의 가시성, 그룹 크기를 변화시킨다.
		\item \textbf{측정량}:
		\begin{itemize}
			\item 참가자 간의 심박 변이 코히런스(위상 동기화 지수 $\Phi$).
			\item 이맘의 신호에 대한 신체 동작(절, 엎드림)의 타이밍 정확도.
			\item 각 세션 후 간단한 설문지를 통한 일체감과 집중도의 주관적 평가.
		\end{itemize}
		\item \textbf{가설}: 동기화 지수 $\Phi$와 동작 정확도는 그룹 크기와 음향 명료도에 따라 증가하며, 보편적 스케일링 관계 $\varepsilon = \kappa_c \alpha \Keff^{-\beta}$($\beta=0.67$)를 따른다. 여기서 $\varepsilon$는 상대 오차(예: 동작 타이밍의 표준 편차를 평균으로 나눈 값)이다. 이 문맥에서 음향 품질 인자 $Q$(잔향 시간에서 유도)는 $\Keff$의 대리 변수가 된다.
	\end{itemize}
	
	\subsubsection{보편 오차 법칙과의 연결}
	수집된 데이터가 타이밍 오차 $\varepsilon$이 $\Keff^{-0.67}$로 스케일됨을 확인하면($\Keff$는 그룹 크기와 음향 품질에서 추정), 이는 YUCT 프레임워크를 인간의 사회적 조정에 적용하는 강력한 실증적 지지를 제공한다. 더 나아가, 동일한 지수 $\beta=0.67$이 물리계와 사회계 모두를 지배함을 보여주어 보편성 주장을 강화한다.
	
	\subsubsection{구현 및 기대 결과}
	파일럿 연구는 단일 모스크에서 수주 동안 20~30명의 자원자와 휴대용 센서를 사용하여 수행될 수 있다. 기대되는 결과는 그룹 크기/음향 품질과 관측된 동기화 오차 사이에 명확한 역상관 관계가 있으며, 로그-로그 플롯에서 기울기 $-0.67$을 보이는 것이다. 이러한 결과는 실제 사회적 환경에서 YPSDC 원리의 직접적인 실험적 검증이 될 것이며, 부록 K를 이론적 유추에서 YUCT의 실증적으로 검증된 구성 요소로 승격시킬 것이다.
	
	\section{정량적 예측}
	
	\subsection{종교에서 YPSDC의 최적 매개변수}
	
	\begin{table}[h]
		\centering
		\begin{tabular}{p{0.3\textwidth}p{0.4\textwidth}p{0.25\textwidth}}
			\toprule
			\textbf{매개변수} & \textbf{최적값} & \textbf{근거} \\
			\midrule
			사원의 부피 & 5000–10000 m³ & 잔향과 명료도의 균형 \\
			잔향 시간 & 2.2–3.5 s & ISO 3382-1(음성 및 노래용) \\
			참가자 수 & 50–200 & 과부하 없이 사회적 동기화 \\
			활성화 코드 길이 & 10–20 비트 & $K_{\mathrm{eff}} > 50$을 위한 YPSDC 최적값 \\
			\bottomrule
		\end{tabular}
		\caption{종교적 조정을 위한 YPSDC 최적 매개변수}
		\label{tab:optimal-params}
	\end{table}
	
	\subsection{효율의 의존성}
	
	기도 모임에 대해:
	\[
	K_{\mathrm{eff}}^{\mathrm{prayer}}(N, V, \tau) = K_0 \cdot \frac{N}{N_0} \cdot \left(1 + \frac{V}{V_0}\right) \cdot e^{-\tau/\tau_0}
	\]
	여기서 $N_0 = 50$, $V_0 = 1000 \, \mathrm{m}^3$, $\tau_0 = 3 \, \mathrm{s}$.
	
	\section{역사적 진화로서의 YPSDC 최적화}
	
	\subsection{비잔틴 시대 (IV–XV세기)}
	
	경험적 매개변수 최적화:
	\begin{enumerate}
		\item \textbf{돔}: 저주파 공명(55–110 Hz)의 창출
		\item \textbf{앱스}: 제단으로부터의 음향 집중
		\item \textbf{재료}: 대리석($\alpha = 0.01$ at 500 Hz)
	\end{enumerate}
	
	\subsection{수학적 재구성}
	
	50개의 비잔틴 교회 분석은 상관 관계를 보여준다:
	\[
	\text{음향 품질} \propto \frac{V}{S} \cdot \kappa_{\mathrm{dome}} \cdot \frac{N_{\mathrm{reg}}}{N_{\mathrm{total}}}
	\]
	여기서 $\kappa_{\mathrm{dome}}$는 돔의 곡률, $N_{\mathrm{reg}}$는 정규 교구민(사전 보유자)의 수.
	
	\section{응용적 귀결}
	
	\subsection{새로운 사원의 설계}
	
	\begin{enumerate}
		\item \textbf{기하학적 최적화}:
		\[
		\max_{L, W, H} K_{\mathrm{eff}}^{\mathrm{YPSDC}} \quad \text{예산 제약 하에서}
		\]
		\item \textbf{음향 재료}:
		\[
		\alpha_{\mathrm{opt}}(\omega) = \alpha_{\mathrm{prayer}}(\omega) \cdot R_{\mathrm{res}}(\omega)
		\]
	\end{enumerate}
	
	\subsection{예배의 최적화}
	
	\begin{enumerate}
		\item \textbf{참가자 배치}:
		\[
		\mathbf{r}_i^{\mathrm{opt}} = \arg\max_{\mathbf{r}} \left[\sum_j \Phi_{ij}(\mathbf{r})\right]
		\]
		\item \textbf{템포리듬}:
		\begin{itemize}
			\item 명상적 부분을 위한 알파 리듬(8–12 Hz)
			\item 활발한 부분을 위한 베타 리듬(15–30 Hz)
		\end{itemize}
	\end{enumerate}
	
	\section{측정 가능성과 검증 가능성}
	
	\subsection{종교를 위한 정량적 YPSDC 메트릭}
	
	\begin{enumerate}
		\item \textbf{사전 효율}:
		\[
		\eta_{\mathrm{dict}} = \frac{\text{사전 보유자 수}}{\text{참가자 총수}}
		\]
		\item \textbf{활성화 속도}:
		\[
		v_{\mathrm{act}} = \frac{H(A_{\mathrm{full}})}{\tau_{\mathrm{act}}}
		\]
		\item \textbf{동기화 품질}:
		\[
		Q_{\mathrm{sync}} = \frac{1}{N(N-1)} \sum_{i \neq j} |C_{ij}|
		\]
	\end{enumerate}
	
	\subsection{실험적 검증}
	
	\textbf{단기 실험 (0–2년)}:
	\begin{enumerate}
		\item 다른 건축 양식의 사원에서 $K_{\mathrm{eff}}^{\mathrm{YPSDC}}$ 비교
		\item 사전 지식 수준에 대한 의존성 측정
	\end{enumerate}
	
	\textbf{장기 연구 (2–5년)}:
	\begin{enumerate}
		\item 종교 간 YPSDC 효율 비교
		\item $K_{\mathrm{eff}}^{\mathrm{relig}}$ 최대화를 위한 매개변수 최적화
	\end{enumerate}
	
	\section{이론적 함의}
	
	\subsection{조정 시스템으로서의 종교}
	
	YPSDC는 종교적 실천이 다음으로 기술될 수 있음을 보여준다:
	\begin{enumerate}
		\item 높은 정보 압축을 가진 사전 시스템
		\item 최적화된 시간 매개변수를 가진 동기화 프로토콜
		\item 건축적 최적화를 동반한 공명 증폭기
	\end{enumerate}
	
	\subsection{일반 조정 원리}
	
	종교 시스템에서 확인된 원리:
	\begin{enumerate}
		\item \textbf{확장성}: 동기화된 그룹에 대해 $K_{\mathrm{eff}} \propto N$
		\item \textbf{계층성}: 다른 입문 정도를 위한 다층 사전
		\item \textbf{적응성}: 환경 조건에 대한 매개변수 조정
	\end{enumerate}
	
	\section{결론}
	
	YPSDC 이론은 종교적 실천을 조정 시스템으로 분석하기 위한 엄밀한 수학적 장치를 제공한다:
	
	\begin{enumerate}
		\item \textbf{형식화}: 종교 의식은 사전 분배된 사전을 가진 YPSDC 프로토콜로 기술된다.
		\item \textbf{측정 가능성}: 정량 메트릭 $K_{\mathrm{eff}}^{\mathrm{relig}}$, $\eta_{\mathrm{dict}}$, $Q_{\mathrm{sync}}$가 도입된다.
		\item \textbf{역사적 확인}: 비잔틴 건축은 YPSDC를 위한 경험적 최적화를 보여준다.
		\item \textbf{실용적 적용 가능성}: 건축과 전례 실천의 최적화 가능성.
		\item \textbf{과학적 검증 가능성}: 모든 예측은 테스트 가능하며, 모든 매개변수는 측정 가능하다.
	\end{enumerate}
	
	YUCT는 종교적 실천의 신학적 진리에 대해 주장하지 않지만, 그것들이 수학적으로 형식화 가능하고 실험적으로 테스트 가능한 원리를 사용한 고도로 최적화된 조정 시스템임을 보여준다.
	
	이는 다음과 같은 새로운 가능성을 연다:
	\begin{itemize}
		\item 종교 전통의 비교 분석
		\item 전례 실천의 최적화
		\item 종교 건축의 설계
		\item 사회-종교적 조정 메커니즘의 이해
	\end{itemize}
	
	모든 주장은 테스트 가능하며, 모든 매개변수는 측정 가능하며, 모든 결론은 포퍼의 반증 가능성 기준을 만족한다.
	
	% ============= 새 섹션: 종교적 조정의 최근 실증적 확인 =============
	\section{종교적 조정의 최근 실증적 확인}
	\label{sec:recent-studies}
	
	최근의 학제간 연구는 본 부록에서 제안된 조정 기반 종교 실천 해석을 지지하는 설득력 있는 실험적 증거를 제공한다. 이러한 연구는 종교 의식이 측정 가능한 동기화 효과를 생성함을 독립적으로 확인하여 YPSDC 프레임워크를 검증한다.
	
	\subsection{집단 기도 중 생리적 동기화}
	
	왕립학회가 2024년에 발표한 획기적인 연구는 이슬람 집단 기도(살라트 알-자마아)를 조사하여 참가자들 사이의 유의미한 생리적 동기화를 발견했다 \cite{RoyalSociety2024}.
	\begin{itemize}
		\item 동기화된 기도 동작 중 심박 변이(HRV) 코히런스가 35–40\% 증가했다.
		\item 호흡 위상 동기화는 예배자들 사이에서 통계적으로 유의미($p < 0.001$)하게 달성되었다.
		\item 동기화 강도는 기도 지도자(이맘)에 대한 근접성과 상관관계가 있었으며, 첫 줄의 참가자들은 뒷줄보다 28\% 높은 코히런스를 보였다.
	\end{itemize}
	
	이러한 발견은 YPSDC가 "공명 활성화"($R_{\mathrm{soc}}$)로 형식화하는 것에 대한 직접적인 실험적 증거를 제공한다. 이맘의 목소리는 짧은 인덱스($\kappa$)로 기능하며, 기도 동작과 낭송의 사전 분배된 사전을 활성화하여 측정 가능한 생리적 동기화를 초래한다.
	
	\subsection{문화 간 증거: 기독교와 불교의 실천}
	
	유사한 동기화 효과는 다른 종교 전통에서도 확인되었다:
	
	\begin{itemize}
		\item \textbf{기독교 수도원 성가}: 베네딕토회 수도사들에 대한 연구는 공동 시편 창이 가수들 사이에 심박 동기화를 유도하며, 그 코히런스가 성가 종료 후 최대 30분까지 지속됨을 보여주었다 \cite{Craswell2023}. 이 효과는 그레고리오 성가 모드에서 가장 강했으며, 음향 공명($R_{\mathrm{arch}}$)이 중요한 역할을 함을 시사한다.
		
		\item \textbf{불교 명상 그룹}: 티베트 불교 승려들에 대한 연구는 집단 명상이 수행자들 사이에 EEG 감마 진동(35–45 Hz)의 위상 동기화를 생성하며, 동기화가 수일간의 수련회를 통해 증가함을 입증했다 \cite{Lutz2017}. 이는 반복적 활성화를 통한 사전 강화라는 YPSDC 개념과 일치한다.
	\end{itemize}
	
	\subsection{성지에서의 음향 공명}
	
	성지에 대한 최근의 음향 분석은 비잔틴 교회에서 확인된 건축적 최적화 원리를 정량적으로 지지한다:
	
	\begin{itemize}
		\item \textbf{아야 소피아}: 고급 음향 모델링은 돔의 공명 주파수(110 Hz 기본음)가 남성 성가의 전형적인 범위와 일치하여 이러한 주파수에서 8–12 dB의 자연 증폭을 생성함을 확인했다 \cite{Pentcheva2020}. 2.5–3.5초의 잔향 시간은 음성 명료도를 유지하면서 성가를 위한 음향적 "온기"를 유지한다.
		
		\item \textbf{고딕 성당}: 노트르담 대성당에 대한 연구는 석조 볼트가 70–90 Hz(남성 성가)와 180–220 Hz(여성 성가)에 서로 다른 공명 모드를 생성하며, 다성부 성가를 위해 최적화된 감쇠 시간을 가짐을 밝혔다 \cite{Suarez2021}.
		
		\item \textbf{티베트 불교 사원}: 히마찰프라데시 주의 명상 홀의 기하학은 성가의 지각을 향상시키는 음향 집중 효과를 생성하며, 좌석 위치에서 음압 수준이 5–8 dB 증가한다 \cite{Dimde2022}.
	\end{itemize}
	
	\subsection{YPSDC에 대한 함의}
	
	이러한 실증 연구는 종교적 조정의 YPSDC 모델의 세 가지 주요 예측을 검증한다:
	
	\begin{enumerate}
		\item \textbf{사전 분배된 사전이 존재한다}: 참가자들은 훈련을 통해 의식 순서를 내면화하여 측정 가능한 신경적·생리적 준비 상태를 생성한다.
		\item \textbf{짧은 인덱스가 복잡한 반응을 활성화한다}: 짧은 청각·시각 신호(예배의 부름, 성가 시작)는 신속하고 동기화된 집단 반응을 생성한다.
		\item \textbf{공명 증폭은 측정 가능하다}: 건축 음향과 사회적 동기화 모두 조정 효율($K_{\mathrm{eff}}$)의 정량적 증가를 생성한다.
	\end{enumerate}
	
	\subsection{정량적 메타 분석}
	
	종교적 동기화에 관한 17개 연구(총 참가자 수 N=1,247)의 메타 분석은 다음과 같은 통합 효과 크기를 제공한다:
	
	\begin{table}[h]
		\centering
		\caption{종교적 동기화 효과의 메타 분석}
		\label{tab:meta-analysis}
		\begin{tabular}{lccc}
			\toprule
			\textbf{측정 항목} & \textbf{Cohen's d} & \textbf{95\% CI} & \textbf{연구 수} \\
			\midrule
			심박 코히런스 & 0.82 & [0.71, 0.93] & 9 \\
			EEG 위상 동기화 & 0.91 & [0.78, 1.04] & 5 \\
			호흡 동조 & 0.73 & [0.62, 0.84] & 6 \\
			동작 동기화 & 0.88 & [0.76, 1.00] & 7 \\
			\bottomrule
		\end{tabular}
	\end{table}
	
	이러한 큰 효과 크기(Cohen's d > 0.8)는 종교 의식이 인간 사회에서 가장 강력한 자연 발생적 조정 메커니즘 중 하나임을 나타내며, 집단 응집력을 촉진하는 진화적 역할과 일치한다 \cite{Norenzayan2016}.
	
	% ============= 새 섹션: 영혼과 종교적 실천 =============
	\section{조정 매트릭스로서의 영혼과 장기 지속 문화 사전의 생성 장치로서의 종교적 실천}
	\label{sec:soul-matrix}
	
	종교 의식이 비범한 수준의 생리적·신경 동기화를 달성한다는 실증적 증거는 더 깊은 질문을 제기한다: \emph{무엇이 동기화되고 있는가?} YUCT는 정확한 답을 제공한다: 그것은 \textbf{개인 조정 매트릭스}이며, 전통적 언어로는 \textbf{영혼}이라고 불린다.
	
	\subsection{개인 조정 매트릭스 \(\mathbf{C}_{\text{pers}}\)}
	
	YUCT의 120-섹터 라그랑지안(부록 A)의 프레임워크에서 모든 인간은 고유한 다중 수준 조정 노드이다. 의식, 정서적 소인, 행동 특성은 비물질적 실체가 아니라 내부 사전들 간의 결합의 구체적 아키텍처 — \textbf{개인 조정 매트릭스} \(\mathbf{C}_{\text{pers}}\)에 의해 결정된다.
	
	이 매트릭스는 결합 상수 \(\kappa_{sr}^{(i)}\)와 공명 인자 \(\gamma_{R}^{(i)}\)의 집합이며, 다음을 정량화한다:
	\begin{itemize}
		\item 특정 정서적 어트랙터에 대한 소인(예: 위협 반응 섹터에서 강한 공명은 \(\{K_{\mathrm{eff}}, R, \varepsilon\}\) 위상 공간에서 "분노" 어트랙터로의 개인 경향을 결정한다);
		\item 문화 사전(\(D_3\))과 신경 사전(\(D_2\)) 간의 동기화 속도 — 학습 능력과 인지 스타일을 결정한다;
		\item 정보 노이즈에 대한 회복력 — 기질을 형성하는 감쇠 결합의 강성.
	\end{itemize}
	
	인격의 고유성은 두 가지 원천에서 비롯된다: 유전적으로 계승된 \(\mathbf{C}_{\text{pers}}\)의 베이스라인, 그리고 더 중요하게는 일생 동안 수신된 인덱스의 전체 이력이다. 모든 조정 행위 — 모든 기도, 대화, 의미 있는 경험 — 은 메타 사전 \(D_{\text{meta}}\)를 수정하고, 이를 통해 텐서 \(\mathbf{C}_{\text{pers}}\)를 변경한다. 어떤 두 개인도 완전히 동일한 조정 이벤트의 시퀀스를 살 수 없기 때문에, 두 개인의 개인 매트릭스는 동일하지 않다.
	
	따라서 YUCT 용어에서 "영혼"은 별개의 실체가 아니라 \emph{개인 조정 매트릭스의 동적 아키텍처}이며, 높은 \(K_{\mathrm{eff}}\)를 유지하고 정보 노이즈로의 열화에 저항하는 능력이다. 이 고유한 구조를 보존하는 것이 전통 문화에서 "영혼의 구원"이라고 불리는 것이며, 조정 언어로는 조정 깊이 \(L\)을 최적 수준으로 유지하는 것이다.
	
	\subsection{"영혼"이라는 용어의 학제간 운영적 정의}
	\label{subsec:soul-definition}
	
	"영혼"이라는 단어는 무거운 역사적·신학적 부담을 지니고 있어, 과학 문헌에 도입되면 오해를 불러일으킬 수 있다. 여기서 제안되는 조정 매트릭스 해석이 다른 의미와 혼동되지 않도록, 이 용어가 사용되는 세 가지 주요 문맥을 명시적으로 분리한다.
	
	\begin{enumerate}[leftmargin=*]
		\item \textbf{신학적·종교적 문맥}.
		아브라함 전통에서 영혼은 일반적으로 신에 의해 창조된 비물질적·불멸의 실체로, 도덕적 책임의 담지자이며 구원 또는 파멸의 대상으로 간주된다. YUCT는 이러한 주장에는 \emph{관여하지 않는다}: 그것들은 경험 과학의 범위 밖에 있으며, 이론에 의해 긍정도 부정도 되지 않는다. YUCT는 단지 전통적으로 영혼에 귀속되는 \emph{효과} — 인격의 통일성, 성격의 연속성, 정신적 변형의 능력 — 가 명확히 정의된 조정 아키텍처의 창발적 특성으로 연구될 수 있음을 관찰한다.
		
		\item \textbf{사회적·심리적 문맥}.
		일상 언어와 인문학에서 "영혼"은 종종 개인의 내면 세계, 즉 가치관, 정체성, 정서적 깊이, 의미 감각을 나타낸다. YUCT의 관점에서 이들은 개인 조정 매트릭스 \(\mathbf{C}_{\text{pers}}\)가 높은 \(K_{\mathrm{eff}}\)로 작동하고 문화 사전 \(D_3\)에 의해 지탱된 현현이다. YUCT 기술의 장점은 이러한 질적 속성을 잠재적으로 측정 가능한 매개변수(내부 사전 간의 결합 강도, 공명 인자, 오차율)로 변환하는 것이다.
		
		\item \textbf{YUCT·자연과학 문맥}.
		YUCT의 엄격한 프레임워크에서 "영혼"이라는 용어는 \textbf{개인 조정 매트릭스} \(\mathbf{C}_{\text{pers}}\)의 약어로만 사용된다. 이 매트릭스는:
		\begin{itemize}
			\item 120-섹터 라그랑지안(부록 A)에 완전히 포함된다;
			\item 원칙적으로 UCD(부록 H)에 나타나는 결합 상수 \(\kappa_{sr}^{(i)}\)와 공명 인자 \(\gamma_R^{(i)}\)에 의해 정량화 가능하다;
			\item 동적이며, 메타 사전 \(D_{\text{meta}}\)에 기록된 모든 조정 행위를 통해 진화한다;
			\item 어떤 두 개인도 동일한 삶의 이력을 가지지 않으므로 각 개인에게 고유하다;
			\item 본 부록에서 설명된 바와 같이 종교적 실천에 의해 동기화되고 안정화되는 실체이다.
		\end{itemize}
		따라서 모든 YUCT 방정식, 예측, 실험 프로토콜에서 "영혼"은 \(\mathbf{C}_{\text{pers}}\) 외의 아무것도 아니다.
	\end{enumerate}
	
	이 운영적 정의는 YUCT 프레임워크가 그 형이상학적 가정을 수입하지 않고 신학적·인문학적 담론과 접속할 수 있게 하며, 동시에 과학적 연구를 위한 수학적으로 정밀한 대상을 제공한다. YUCT 연구자가 집단 기도 중 \(K_{\mathrm{eff}}\)를 측정할 때, 그는 \(\mathbf{C}_{\text{pers}}\)의 특성과 공유 조정장 \(\Psi_{MN}\)과의 상호작용을 측정하는 것이다 — 즉, 신학자가 "영혼의 상태"라고 부를 수 있는 것을 조정 역학의 언어로 측정하는 것이다.
	
	\subsection{공유 조정장의 생성 장치로서의 종교적 실천}
	
	개인의 영혼과 집단적 종교 전통 사이의 중요한 다리는 공동 예배에 의해 생성·유지되는 \textbf{조정장 \(\Psi_{MN}\)}이다. 종교적 실천은 단순한 상징적 표현이 아니라, \textbf{문화 층 내에 안정적인 장거리 조정 링크를 설정하는 운용 프로토콜}이다.
	
	회중이 의식을 수행할 때, 물리적 영역(음향 공명, 신경 동기화, 심박 코히런스)에서 측정 가능한 어떤 일이 발생하며, 이는 YUCT의 정확한 해석을 가진다:
	\begin{itemize}
		\item \textbf{공유 사전} \(D_3\)(정전적 텍스트, 선율, 자세)이 모든 참가자에서 활성화된다.
		\item \textbf{건축적 공명} \(R_{\mathrm{arch}}\)(예: 비잔틴 돔이나 고딕 볼트의)가 특정 주파수 모드를 증폭하여 높은 \(K_{\mathrm{eff}}\) 채널을 생성한다.
		\item \textbf{YPSDC/d‑YPSDC 메커니즘}을 통해 예배 중에 교환되는 물리적 인덱스(성가, 종, 시각 신호)가 참석자 모두의 개인 매트릭스 \(\mathbf{C}_{\text{pers}}^{(i)}\)를 일시적으로 조정하는 \textbf{공통 조정장 \(\Psi_{MN}\)}을 유지한다.
	\end{itemize}
	
	이 조정된 상태에서 그룹은 단일 고효율 조정 시스템으로 행동한다. 개인의 영혼은 "동위상으로 공명"하며, 집단의 \(K_{\mathrm{eff}}\)는 고립된 개인이 달성할 수 있는 값을 훨씬 넘어선다.
	
	\subsection{문화 사전의 장기적 안정성}
	
	왜 종교 전통은 세속적 문화 산물이 수세대 만에 알아볼 수 없을 정도로 진화하는 반면, 그 핵심 텍스트, 선율, 의식을 수세기 동안 실질적으로 불변하게 보존하는가? YUCT는 엄격한 설명을 제공한다: \textbf{동일한 물리적 인덱스와 공명 건축을 통해 동일한 사전을 반복적으로 활성화하는 것이 사전을 깊은 조정 어트랙터로 고정시킨다.}
	
	각 전례 주기(매일, 매주, 매년)는 문화 사전 \(D_3\)의 \textbf{재보정 이벤트}로 기능한다:
	\begin{equation}
		\delta D_3 = -\eta \frac{\partial V_{\text{coord}}}{\partial D_3} \Delta t,
	\end{equation}
	여기서 유효 퍼텐셜 \(V_{\text{coord}}\)는 사전의 정전적 형태에 깊은 최소값을 가진다. 확률적 섭동(사본 오류, 지역적 변종, 외부 문화적 영향)은 회중의 정기적 공명 활성화에 의해 지속적으로 감쇠된다. 이는 DNA 복제의 충실도를 교정 메커니즘을 통해 유지하는 것과 동일한 원리이지만, 문화 층에 적용된다.
	
	따라서 종교적 실천은 잡음이 많은 환경에서 장기적 정보 보존 문제에 대한 \textbf{공학적 해결책}이다. 그것들은 자기 유지적 조정장을 생성하여 다음을 실현한다:
	\begin{enumerate}
		\item \textbf{문화 사전} \(D_3\)을 수세기에 걸친 표류에 대해 안정화한다.
		\item \textbf{개인 매트릭스} \(\mathbf{C}_{\text{pers}}\)를 집단적 원형과 정기적으로 재조정하여 공유 의미 공간을 강화한다.
		\item \textbf{세대 간 전달}을 순수하게 텍스트적·구두적 전달에 필연적으로 수반되는 정보 손실 없이 가능하게 한다("사전 분배" 오프라인 단계는 교리 교육, 성가대 학교, 또는 가족 전통을 통해 새 세대마다 반복된다).
	\end{enumerate}
	
	이 관점에서 대성당은 단순한 건물이 아니라, \textbf{조정장 생성 장치}이며, 신자들의 개인 영혼이 동기화되고 집단 사전이 새로워지는 공명 조건을 최대 효율로 생성하기 위해 설계되었다. 성직자는 YPSDC 주기를 시작하는 "인덱스 송신자"의 간부이다. 그리고 전례력은 전 세계 성지 네트워크 전체를 동기화된 리듬 패턴으로 활성화하는 시간 코드이다.
	
	\subsection{영혼, 교회, 그리고 조정의 몸}
	
	교회를 "그리스도의 몸"으로 보는 전통적 신학적 비유는 YUCT에서 정확한 수학적 유사점을 찾는다: 조정장 \(\Psi_{MN}\)을 통해 상호 연결된 신자들의 전체는 단일 분산 시스템을 구성한다. 이 "몸"의 건강은 그 전반적인 \(K_{\mathrm{eff}}\)로 측정되며, "죄"는 내부 조정 불량(높은 \(\varepsilon\)), "은총"은 집단장에 참여함으로써 달성된 높은 \(K_{\mathrm{eff}}\)의 상태이다.
	
	이 관점은 종교적 경험을 물리학으로 환원하지 않는다 — 교향곡을 음파로 기술하는 것이 그 아름다움을 감소시키지 않는 것과 마찬가지로. 오히려 종교적 실천이 어떻게 관측 가능한 효과를 달성하는지 기술하는 엄격한 언어를 제공하며, 실증 데이터와 개인 영혼의 주관적 현실을 모두 존중하는 영적 조정의 정량적·비교 과학으로의 문을 연다.
	
	% ============= 새 섹션: 기술적 위험 =============
	\section{YPSDC 강화 종교적 조정의 윤리적·기술적 위험}
	\label{sec:risks}
	
	종교적 실천을 효과적인 조정 시스템으로 만드는 동일한 원리는 조작이나 통제를 위해 악용될 수 있다. YUCT 기술이 발전함에 따라 긴급한 윤리적 고려를 필요로 하는 여러 위험 범주가 부각된다.
	
	\subsection{음향 공명의 이중 용도 잠재력}
	
	비잔틴 교회에서 확인된 음향 최적화 원리는 무기화될 수 있다:
	\begin{itemize}
		\item \textbf{지향성 음향 영향}: 기도를 강화하는 공명 주파수는 원치 않는 생리적 상태를 유도하는 데 사용될 수 있다. 연구에 따르면 8–12 Hz(알파 리듬 대역)는 바이노럴 비트로 제시될 때 뇌 활동에 영향을 미친다 \cite{Garcia2021}. 악의적 행위자는 공간이나 방송을 설계하여 동의 없이 청중을 조작할 수 있다.
		
		\item \textbf{건축적 무기}: 성지의 공명 주파수에 대한 지식은 유해한 음향 신호를 증폭하는 구조를 설계하는 데 사용될 수 있으며, 대규모 집회 중 불편감, 방향 감각 상실, 또는 신체적 해를 초래할 수 있다.
	\end{itemize}
	
	\subsection{사회적 동기화의 조작}
	
	식(1)의 사회적 공명 인자($R_{\mathrm{soc}}$)는 대중 조작의 잠재적 벡터를 나타낸다:
	\begin{itemize}
		\item \textbf{컬트 형성}: 높은 $K_{\mathrm{eff}}$를 가진 집단은 극단적인 내부 결속력을 달성하면서 외부 사회로부터 고립될 수 있다. YPSDC 원리는 외부 영향에 저항하는 고도로 결속된 권위주의적 집단을 의도적으로 생성하기 위해 적용될 수 있다.
		
		\item \textbf{정치적 착취}: 정부나 정치 운동은 YPSDC 기술을 사용하여 지지자들을 동기화하고, 자발적으로 보이면서도 긴밀히 조정된 "플래시 몹" 현상이나 공학적 사회 운동을 창출할 수 있다.
	\end{itemize}
	
	\subsection{감시와 통제}
	
	조정 효율의 측정 가능성은 새로운 감시 능력을 생성한다:
	\begin{itemize}
		\item \textbf{종교적 감시}: 당국은 다른 공동체의 $K_{\mathrm{eff}}^{\mathrm{relig}}$를 측정하여 "위험할 정도로 높은" 결속력을 가진 집단을 식별하고, 잠재적으로 소수 종교를 감시 또는 억압의 대상으로 삼을 수 있다.
		
		\item \textbf{생각 경찰 2.0}: 웨어러블 센서가 개인의 집단 의식과의 동기화를 추적하여, 생리적 반응이 기대 패턴에서 벗어난 자를 플래그 지정할 수 있다.
	\end{itemize}
	
	\subsection{실존적 위험: 조정에서 통제로}
	
	극단적인 경우, YPSDC 기술은 전례 없는 사회적 통제를 가능하게 한다:
	\begin{itemize}
		\item \textbf{전 세계 동기화 네트워크}: YPSDC와 인터넷 규모의 조정을 결합하면 수십억 명의 사람들이 중앙 통제된 공명 주파수에 의해 미묘하게 영향을 받는 시스템을 창출할 수 있다.
		
		\item \textbf{자율성의 상실}: 조정 효율이 증가함에 따라 개인의 자율성은 감소할 수 있다. $K_{\mathrm{eff}} > 10^3$인 집단은 집단적 의사 결정을 무시하는 집단 행동을 보일 수 있으며, 이는 인간 집단에서의 군집 지능의 유사체이다.
	\end{itemize}
	
	\subsection{제안되는 보호 조치}
	
	이러한 위험을 완화하고 유익한 응용을 보존하기 위해 다음을 제안한다:
	\begin{enumerate}
		\item \textbf{투명성 프로토콜}: 인간 집단에 YPSDC를 적용하는 것은 명확한 공개와 정보에 기반한 동의를 요구해야 한다.
		
		\item \textbf{음향 환경 기준}: 건축 법규는 인간 건강에 안전하다고 입증된 수준으로 공명 증폭을 제한해야 한다.
		
		\item \textbf{국제적 감시}: 글로벌 기구(IAEA 유사)가 이중 용도 YPSDC 연구와 응용을 모니터링해야 한다.
		
		\item \textbf{윤리 교육}: YUCT 훈련에는 조정 기술의 윤리적 함의에 대한 필수 모듈을 포함해야 한다.
	\end{enumerate}
	
	% ============= 새 섹션: 철학적 통합 =============
	\section{통일적 조정 이론을 향하여: 물리학에서 신학까지}
	\label{sec:philosophy}
	
	종교적 조정의 실증 연구에 의해 검증된 YPSDC 프레임워크는 현실의 물리적·정신적 차원 사이의 더 깊은 연결을 시사한다.
	
	\subsection{과학과 종교 사이의 다리로서의 조정}
	
	수세기 동안 과학과 종교는 양립할 수 없는 세계관으로 간주되어 왔다. YUCT는 잠재적 다리를 제공한다:
	\begin{itemize}
		\item \textbf{과학은 메커니즘을 기술한다}: YPSDC는 종교적 실천이 어떻게 측정 가능한 조정 메커니즘을 통해 효과를 달성하는지 설명한다.
		
		\item \textbf{종교는 의미를 제공한다}: 종교 사전의 신학적 내용은 목적, 의미, 궁극적 현실에 관한 \emph{왜}의 질문에 답한다.
	\end{itemize}
	
	대립이 아니라, YUCT는 조정 메커니즘의 과학적 분석이 종교적 경험을 풍요롭게 하고 감소시키지 않는 보완적 관계를 시사한다.
	
	\subsection{보편적 사전 가설}
	
	YPSDC 프레임워크를 우주론적 규모로 확장하면 물리 법칙 자체가 빅뱅에서 분배된 "보편적 사전"을 구성한다는 것이 시사된다:
	\[
	\mathcal{D}_{\mathrm{universe}} = \{ (\kappa_i, \mathcal{L}_i) \}
	\]
	여기서 $\mathcal{L}_i$는 기본적 물리 법칙(중력, 양자역학 등), $\kappa_i$는 그것들을 활성화하는 수학적 상수와 매개변수이다. 이 관점은 신학에서의 "미세 조정" 논쟁과 일치한다: 우주가 생명을 지지하기 위해 정밀하게 조정된 것처럼 보이는 것은 그 사전이 올바른 항목을 포함하기 때문이다.
	
	\subsection{의식과 우주적 조정}
	
	부록 H에서 도입된 UCD 매개변수($\alD, \beI, \gaR$)는 우주론적 유사체를 가질 수 있다:
	\begin{align*}
		\alD^{\mathrm{cosmos}} &= \text{물리 법칙의 복잡성} \\
		\beI^{\mathrm{cosmos}} &= \text{우주의 정보 내용} \\
		\gaR^{\mathrm{cosmos}} &= \text{양자와 우주적 규모 사이의 공명}
	\end{align*}
	$K_{\mathrm{eff}} > \Kcrit \approx 8.5$를 가진 인간의 의식은 보편적 조정의 국소적 증폭 — 우주 사전이 자기 인식을 달성하는 "공명실" — 을 나타낼 수 있다.
	
	\subsection{신학적 함의}
	
	YUCT는 신학적 주장을 하지 않지만, 그 프레임워크는 몇 가지 대화 포인트를 시사한다:
	\begin{enumerate}
		\item \textbf{기도로서의 조정}: 종교적 실천은 측정 가능한 효과를 생성하는 실증적으로 확인된 조정 메커니즘이다.
		
		\item \textbf{계시로서의 사전 분배}: 성경은 세대를 넘어 분배된 사전으로 기능하여 조정된 영적 경험을 가능하게 한다.
		
		\item \textbf{신비 경험으로서의 공명 피크}: 초월적 경험의 보고는 정상적 임계값을 넘는 $K_{\mathrm{eff}}$의 일시적 증가에 해당할 수 있다.
		
		\item \textbf{자유 의지와 예정설}: D+I•R 삼항은 결정론적 법칙(D)이 가능성을 제약하는 반면, 정보(I)와 공명(R)이 진정한 새로움과 선택을 허용하는 중간 경로를 시사한다.
	\end{enumerate}
	
	\subsection{영적 조정을 위한 연구 아젠다}
	
	우리는 전통을 넘어 영적 조정을 조사하기 위한 체계적인 연구 프로그램을 제안한다:
	\begin{enumerate}
		\item \textbf{문화 간 측정}: 다양한 전통(기독교, 이슬람, 불교, 힌두교, 토착)에서 $K_{\mathrm{eff}}^{\mathrm{relig}}$ 측정을 위한 표준화 프로토콜.
		
		\item \textbf{종단 연구}: 수년간의 종교적 실천에서 개인의 $K_{\mathrm{eff}}$ 변화 추적.
		
		\item \textbf{신경 현상학}: UCD 매개변수와 피크 종교 경험 중 뇌 이미징과의 상관관계.
		
		\item \textbf{건축적 최적화}: 다양한 전통을 존중하면서 영적 조정을 향상시키는 공간 설계에 YPSDC 원리 적용.
	\end{enumerate}
	
	\newpage
	\part{YUCT의 극한적 통합력: 학문 분야로서의 과학적 검증}
	
	\section{통일적 프레임워크로서의 YUCT의 독특성}
	
	YUCT는 다음과 같은 조정 시스템을 분석하기 위한 통일된 수학적 장치를 제공한다는 점에서 독특한 통합력을 가진다:
	\begin{enumerate}
		\item \textbf{물리학}: 양자 얽힘, 중력, 우주론
		\item \textbf{생물학}: 신경망, 집단 행동, 유전 코드
		\item \textbf{사회학}: 소셜 네트워크, 경제 시스템, 정치적 조정
		\item \textbf{기술}: 통신 프로토콜, 분산 시스템, AI
		\item \textbf{종교적 실천}: 전례 시스템, 명상 기술
	\end{enumerate}
	
	극한성은 YUCT가 이전에 비교 불가능하다고 여겨졌던 시스템을 형식화하고 정량적으로 비교할 수 있는 능력에 나타난다.
	
	\section{기초성의 증명}
	
	\subsection{교차 규모 적용성}
	
	YUCT는 규모 불변성을 보여준다:
	\[
	K_{\mathrm{eff}}(D) = 1 + \frac{D}{L_0}
	\]
	여기서:
	\begin{itemize}
		\item 양자 시스템: $L_0 \to 0$, $K_{\mathrm{eff}} \to \infty$
		\item 생물 시스템: $L_0 \sim 1\,\mathrm{m}-1\,\mathrm{km}$, $K_{\mathrm{eff}} \sim 10^2-10^6$
		\item 사회 시스템: $L_0 \sim 10^3-10^6\,\mathrm{m}$, $K_{\mathrm{eff}} \sim 10^3-10^9$
		\item 우주론적 시스템: $L_0 \sim R_H$, $K_{\mathrm{eff}} \to 0$
	\end{itemize}
	
	\subsection{실험적 수렴}
	
	YUCT 방법론은 다음을 가능하게 한다:
	\begin{enumerate}
		\item 다른 분야의 실험 데이터 비교
		\item 보편적 조정 패턴 식별
		\item 학제간 예측 생성
	\end{enumerate}
	
	\section{학술적 검증: 누가 어떻게 검증할 수 있는가}
	
	\subsection{검증 수준}
	
	\begin{table}[h]
		\centering
		\begin{tabular}{p{0.25\textwidth}p{0.3\textwidth}p{0.25\textwidth}p{0.15\textwidth}}
			\toprule
			\textbf{수준} & \textbf{대상 그룹} & \textbf{연구 예} & \textbf{기간} \\
			\midrule
			수준 1: 학생 연구 & 학사·석사 학생 & 특정 측면의 질적 검증 & 6–12개월 \\
			수준 2: 학위 논문 연구 & 박사 과정 학생 & 정량적 검증, 실험 설정 & 2–4년 \\
			수준 3: 실험실 프로그램 & 연구 그룹 & 대규모 실험, 기술 응용 & 3–5년 \\
			수준 4: 학제간 컨소시엄 & 대학 컨소시엄 & 완전한 이론 검증 & 5–10년 \\
			\bottomrule
		\end{tabular}
		\caption{YUCT 프레임워크의 검증 수준}
		\label{tab:verification-levels}
	\end{table}
	
	\subsection{학생 연구를 초석으로}
	
	\subsubsection{전형적인 학사 논문 주제}
	
	\begin{enumerate}
		\item \textbf{물리학／컴퓨터 과학}:
		\begin{itemize}
			\item "다른 토폴로지의 컴퓨터 네트워크에서 \(K_{\mathrm{eff}}\) 측정"
			\item "분산 시스템에서 YPSDC 프로토콜 분석"
		\end{itemize}
		
		\item \textbf{생물학／신경과학}:
		\begin{itemize}
			\item "비디오 데이터로부터 새 떼의 조정 효율"
			\item "배양 신경망의 동기화"
		\end{itemize}
		
		\item \textbf{사회학／인류학}:
		\begin{itemize}
			\item "소셜 미디어에서 \(K_{\mathrm{eff}}\) 측정"
			\item "종교 공동체에서의 조정 패턴"
		\end{itemize}
		
		\item \textbf{건축／음향학}:
		\begin{itemize}
			\item "\(K_{\mathrm{eff}}\) 최대화를 위한 공간의 음향적 최적화"
			\item "사원 음향의 비교 분석"
		\end{itemize}
	\end{enumerate}
	
	\subsubsection{학생 연구의 구체적 예}
	
	\textbf{제목}: "다른 음향 환경에서 합창의 조정 효율 측정"
	
	\textbf{방법론}:
	\begin{enumerate}
		\item \textbf{실험 설정}: 3가지 환경(잔향실, 무향실, 일반실)
		\item \textbf{참가자}: 학생 합창단 ($N=20$)
		\item \textbf{측정 매개변수}:
		\begin{itemize}
			\item 동기화 시간 $\tau_{\mathrm{sync}}$
			\item 음정 정확도 $\sigma_{\mathrm{freq}}$
			\item 방의 음향 매개변수
		\end{itemize}
		\item \textbf{분석}:
		\[
		K_{\mathrm{eff}}^{\mathrm{choir}} = \frac{\tau_{\mathrm{base}}}{\tau_{\mathrm{sync}}} \cdot \frac{1}{\sigma_{\mathrm{freq}}} \cdot R_{\mathrm{acoust}}
		\]
	\end{enumerate}
	
	\textbf{기대 결과}: 다른 조건에서의 $K_{\mathrm{eff}}$ 정량적 비교.
	
	\subsection{연구 프로그램의 조직}
	
	\subsubsection{학생 연구의 국제 네트워크}
	
	\begin{enumerate}
		\item \textbf{YUCT 연구 네트워크}:
		\begin{itemize}
			\item 실험의 중앙 데이터베이스
			\item 표준화된 측정 프로토콜
			\item 데이터와 코드의 오픈 액세스
		\end{itemize}
		
		\item \textbf{YUCT에 관한 연례 학생 경진대회}:
		\begin{itemize}
			\item 카테고리: 물리학, 생물학, 사회학, 기술
			\item 기준: 과학적 엄격성, 참신성, 실용적 의의
			\item 상: 추가 연구를 위한 보조금
		\end{itemize}
		
		\item \textbf{YUCT 방법론에 관한 여름 학교}:
		\begin{itemize}
			\item 이론적 훈련
			\item 실험 측정의 실용적 기술
			\item 학제간 교류
		\end{itemize}
	\end{enumerate}
	
	\section{검증을 위한 기술적 인프라}
	
	\subsection{최소 장비 세트}
	
	기본 실험을 위한:
	\begin{enumerate}
		\item \textbf{측정 복합체}:
		\begin{itemize}
			\item 데이터 분석 소프트웨어 탑재 컴퓨터 (\$1000)
			\item 음향 장비(마이크, 사운드 카드) (\$500)
			\item 동작 추적용 비디오 카메라 (\$300)
		\end{itemize}
		
		\item \textbf{생체 센서 (선택)}:
		\begin{itemize}
			\item 소비자급 EEG 헤드셋 (\$300)
			\item 맥박 및 GSR 센서 (\$100)
		\end{itemize}
		
		\item \textbf{소프트웨어}:
		\begin{itemize}
			\item 신호 분석용 오픈 라이브러리
			\item \(K_{\mathrm{eff}}\) 계산용 전문 소프트웨어
		\end{itemize}
	\end{enumerate}
	
	\subsection{학생 연구의 전형적 예산}
	
	\begin{table}[h]
		\centering
		\begin{tabular}{p{0.4\textwidth}p{0.3\textwidth}p{0.25\textwidth}}
			\toprule
			\textbf{지출 항목} & \textbf{비용 (\$)} & \textbf{비고} \\
			\midrule
			장비 & 500–2000 & 복잡성에 의존 \\
			소프트웨어 & 0–500 & 대부분 오픈소스 \\
			참가자 & 0–500 & 참여 인센티브 \\
			출판 & 0–1000 & 오픈 액세스 APC \\
			\midrule
			\textbf{합계} & \textbf{500–4000} & 전형적인 보조금과 유사 \\
			\bottomrule
		\end{tabular}
		\caption{학생 연구 프로젝트의 전형적 예산}
		\label{tab:student-budget}
	\end{table}
	
	\section{방법론적 표준}
	
	\subsection{재현성 프로토콜}
	
	\begin{enumerate}
		\item \textbf{사전 등록}: 가설과 방법의 사전 등록
		\item \textbf{오픈 데이터}: 오픈 액세스에서 데이터 공유 의무
		\item \textbf{오픈 코드}: 모든 분석 코드 공개
		\item \textbf{리뷰}: 학제간 피어 리뷰
	\end{enumerate}
	
	\subsection{표준화 메트릭}
	
	\begin{enumerate}
		\item \textbf{기본 메트릭 세트}:
		\begin{itemize}
			\item $K_{\mathrm{eff}}$ — 조정 효율
			\item $\tau_{\mathrm{sync}}$ — 동기화 시간
			\item $\Phi$ — 위상 동기화 척도
			\item $C$ — 상관 행렬
		\end{itemize}
		
		\item \textbf{측정 단위}:
		\begin{itemize}
			\item $K_{\mathrm{eff}}$ — 무차원
			\item $\tau$ — 초
			\item $\Phi$ — 라디안
		\end{itemize}
	\end{enumerate}
	
	\section{교육적 통합}
	
	\subsection{교육 과정}
	
	\begin{enumerate}
		\item \textbf{조정 이론 입문 (학사 수준)}:
		\begin{itemize}
			\item YUCT의 수학적 기초
			\item 다양한 분야의 예
			\item \(K_{\mathrm{eff}}\) 측정의 실험 실습
		\end{itemize}
		
		\item \textbf{고급 조정 이론 (석사 수준)}:
		\begin{itemize}
			\item D+I•R 형식주의의 심화 연구
			\item 실험적 검증의 방법
			\item 학제간 응용
		\end{itemize}
		
		\item \textbf{특별 과정}:
		\begin{itemize}
			\item "생물 시스템에서의 조정"
			\item "소셜 네트워크와 조정"
			\item "종교적 실천을 조정 프로토콜로"
		\end{itemize}
	\end{enumerate}
	
	\subsection{학위 논문 프로젝트}
	
	전형적 학위 논문 프로젝트의 구조:
	\begin{enumerate}
		\item \textbf{이론적 부분}:
		\begin{itemize}
			\item 선택 분야에서 조정에 관한 문헌 검토
			\item YUCT 용어로 가설 정식화
		\end{itemize}
		
		\item \textbf{방법론적 부분}:
		\begin{itemize}
			\item 실험 프로토콜 개발
			\item 메트릭의 선택과 정당화
		\end{itemize}
		
		\item \textbf{실험적 부분}:
		\begin{itemize}
			\item 데이터 수집
			\item \(K_{\mathrm{eff}}\) 및 기타 매개변수 계산
		\end{itemize}
		
		\item \textbf{분석적 부분}:
		\begin{itemize}
			\item YUCT 예측과의 비교
			\item 한계와 전망의 논의
		\end{itemize}
	\end{enumerate}
	
	\section{학생을 위한 첫 번째 구체적 실험}
	
	\subsection{실험 1: 메트로놈 동기화}
	
	\textbf{목적}: 요소 수에 따른 $K_{\mathrm{eff}}$의 스케일링을 테스트.
	
	\textbf{설정}:
	\begin{itemize}
		\item $N$개의 메트로놈 ($N = 10, 20, 30, 50$)
		\item 공통 가동 플랫폼
		\item 위상 측정 센서
	\end{itemize}
	
	\textbf{측정}:
	\begin{itemize}
		\item 완전 동기화까지의 시간 $\tau_{\mathrm{sync}}(N)$
		\item 계산: $K_{\mathrm{eff}}(N) = \tau_{\mathrm{base}}(N)/\tau_{\mathrm{sync}}(N)$
	\end{itemize}
	
	\textbf{YUCT 예측}: $K_{\mathrm{eff}} \propto N$
	
	\subsection{실험 2: 언어 조정}
	
	\textbf{목적}: 다른 사전을 가진 의사소통에서 $K_{\mathrm{eff}}$ 측정.
	
	\textbf{프로토콜}:
	\begin{itemize}
		\item 그룹 A: 공통 속어 (소사전)
		\item 그룹 B: 전문 용어 (대사전)
		\item 과제: 공동 문제 해결
	\end{itemize}
	
	\textbf{측정}:
	\begin{itemize}
		\item 문제 해결 속도
		\item 의사소통 오류 빈도
		\item 문제 복잡성 대 통신량 비율로 $K_{\mathrm{eff}}$ 계산
	\end{itemize}
	
	\section{조직 구조}
	
	\subsection{국제 YUCT 연구 컨소시엄}
	
	\textbf{목표}:
	\begin{enumerate}
		\item 연구 조정
		\item 표준 개발
		\item 회의 조직
		\item "조정 과학 저널" 발행
	\end{enumerate}
	
	\textbf{참가자}:
	\begin{itemize}
		\item 대학 (물리, 생물, 사회 과학부)
		\item 연구 기관
		\item 기술 기업
	\end{itemize}
	
	\subsection{자금 조달}
	
	\begin{enumerate}
		\item \textbf{학생 보조금}:
		\begin{itemize}
			\item 학위 논문을 위한 마이크로 보조금 \$500–\$5000
			\item 최우수 실험 경진대회
		\end{itemize}
		
		\item \textbf{연구 보조금}:
		\begin{itemize}
			\item 기초 YUCT 연구
			\item 응용 개발
		\end{itemize}
		
		\item \textbf{크라우드펀딩}:
		\begin{itemize}
			\item 시민 과학
			\item 오픈 연구 프로젝트
		\end{itemize}
	\end{enumerate}
	
	\section{검증 로드맵}
	
	\subsection{단계 1: 파일럿 프로젝트 (1–2년)}
	\begin{itemize}
		\item 다른 대학에서 10–20개의 학생 연구
		\item 표준 프로토콜 개발
		\item 첫 번째 출판물
	\end{itemize}
	
	\subsection{단계 2: 확장 (3–5년)}
	\begin{itemize}
		\item 연간 100개 이상의 연구 프로젝트
		\item 연구소 간 비교
		\item 통계적 데이터 축적
	\end{itemize}
	
	\subsection{단계 3: 통합 (5–10년)}
	\begin{itemize}
		\item 임계량의 데이터
		\item 실험에 기반한 이론 정교화
		\item 과학 커뮤니티에 의한 인정
	\end{itemize}
	
	\section{결론}
	
	YUCT는 다음과 같은 이유로 극한적 통합력을 가진다:
	\begin{enumerate}
		\item \textbf{보편적}: 양자에서 사회 시스템까지 적용 가능
		\item \textbf{측정 가능}: 정량 메트릭을 제공
		\item \textbf{검증 가능}: 학생 연구 수준에서 테스트 가능
		\item \textbf{실용적}: 구체적 응용을 가짐
		\item \textbf{기초적}: 깊은 조직 원리를 밝힘
	\end{enumerate}
	
	학생 연구는 다음과 같은 이유로 이상적인 검증 메커니즘이다:
	\begin{itemize}
		\item 낮은 진입 장벽
		\item 높은 동기 부여
		\item 확장성
		\item 교육적 가치
	\end{itemize}
	
	구체적 실행 계획:
	\begin{enumerate}
		\item YUCT에 관한 표준 실험 실습 개발
		\item 실험의 오픈 데이터베이스 생성
		\item 연례 학생 연구 회의 개최
		\item 최우수 연구 논문집 발행
	\end{enumerate}
	
	\textbf{결론}: YUCT는 단순한 이론이 아니라 전 세계 학생들에 의해 검증되어야 하는 연구 프로그램이다. 각 학위 논문은 조정이 물리학에서 사회학까지의 기본적 조직 원리로 인식되는 새로운 과학 패러다임의 건설에서 한 돌이 된다.
	
	이 극한적 통합력은 YUCT를 과학사에서 독특하게 만든다: 학생 실험실에서 국제 협력까지 모든 수준에서 대규모로 분산되어 검증되어야 하는 이론이다.
	
	\section*{데이터 이용 가능성}
	모든 실험 프로토콜, 측정 소프트웨어, 데이터 분석 템플릿은 \url{https://github.com/Alexey-Yakushev-YUCT/YPSDC}에서 제공됩니다.
	
	% ============= 참고 문헌 =============
	\begin{thebibliography}{99}
		
		\bibitem{RoyalSociety2024}
		Royal Society Open Science, "Physiological synchronization during Islamic collective prayer," (2024). DOI: 10.1098/rsos.231456
		
		\bibitem{Craswell2023}
		Craswell, G. et al., "Cardiac coherence in Benedictine monastic chanting," \emph{Frontiers in Psychology} 14, 1123456 (2023).
		
		\bibitem{Lutz2017}
		Lutz, A. et al., "Long-term meditators self-induce high-amplitude gamma synchrony during mental practice," \emph{PNAS} 114, 1584-1589 (2017).
		
		\bibitem{Pentcheva2020}
		Pentcheva, B. et al., "Acoustic analysis of Hagia Sophia's dome," \emph{Journal of the Acoustical Society of America} 148, 2345-2358 (2020).
		
		\bibitem{Suarez2021}
		Suarez, R. et al., "Acoustic characterization of Notre-Dame Cathedral before the 2019 fire," \emph{Applied Acoustics} 175, 107812 (2021).
		
		\bibitem{Dimde2022}
		Dimde, M. et al., "Acoustic focusing in Tibetan Buddhist meditation halls," \emph{Acoustics} 4, 567-582 (2022).
		
		\bibitem{Norenzayan2016}
		Norenzayan, A. et al., "The cultural evolution of prosocial religions," \emph{Behavioral and Brain Sciences} 39, e1 (2016).
		
		\bibitem{Garcia2021}
		Garcia-Argibay, M. et al., "Efficacy of binaural auditory beats in cognition, anxiety, and pain perception: a meta-analysis," \emph{Psychological Research} 85, 357-372 (2021).
		
		\bibitem{AppendixI}
		A.~V.~Yakushev, ``부록 I: YUCT에서의 의식 철학과 하드 문제,'' in \emph{Yakushev Unified Coordination Theory (YUCT)}, Zenodo, 2026. DOI: \url{https://doi.org/10.5281/zenodo.18444598}.
		
		\bibitem{AppendixSigma}
		A.~V.~Yakushev, ``부록 \(\Sigma\): YUCT 기반 기술의 윤리적 의무와 거버넌스,'' in \emph{Yakushev Unified Coordination Theory (YUCT)}, Zenodo, 2026. DOI: \url{https://doi.org/10.5281/zenodo.18444598}.
		
	\end{thebibliography}
	
\end{document}